% Generated draft: requires mathematical and source review.
\documentclass[11pt]{article}
\usepackage[margin=25mm]{geometry}
\usepackage{fontspec}
\setmainfont{Latin Modern Roman}
\usepackage{amsmath,amssymb,mathtools,unicode-math}
\setmathfont{Latin Modern Math}
\usepackage{enumitem,needspace,xurl,hyperref}
\setlength{\parindent}{0pt}
\setlength{\parskip}{5pt}
\setlength{\emergencystretch}{3em}
\newcommand{\N}{\mathbb{N}}
\newcommand{\Z}{\mathbb{Z}}
\newcommand{\Q}{\mathbb{Q}}
\newcommand{\R}{\mathbb{R}}
\newcommand{\C}{\mathbb{C}}
\newcommand{\F}{\mathbb{F}}
\newcommand{\A}{\mathbb{A}}
\newcommand{\PP}{\mathbb{P}}
\newcommand{\E}{\mathbb{E}}
\newcommand{\eps}{\varepsilon}
\newcommand{\dd}{\,\mathrm{d}}
\begin{document}
\section*{UnsolvedMath ID: 30000839}
\textbf{Source problem code:} catalogue entry
\subsection{Definitions and mathematical statement}
In the standard continuous-time one-dimensional TASEP on $\Z$, a configuration is $\eta\in\{0,1\}^{\Z}$, where $\eta(i)=1$ indicates a particle. Each particle attempts to jump from $i$ to $i+1$ at rate one, and the jump occurs only when the destination is empty. Equivalently, on cylinder functions the generator is \[(\mathcal LF)(\eta)=\sum_{i\in\Z}\eta(i)(1-\eta(i+1))\bigl(F(\eta^{i,i+1})-F(\eta)\bigr),\] where $\eta^{i,i+1}$ interchanges the two indicated occupations. This is the usual model interpretation; the extract does not itself specify state space, time convention, or rates. Its question is: for the various classes of initial conditions, what are the limiting one-point fluctuation laws and the limiting fluctuation process as $t\to\infty$? To formulate a particular limit one must choose an observable $O_t$, observation locations $x_t(s)$, deterministic centering $a_t(s)$, and positive fluctuation scale $b_t(s)$, and ask for limits of \[Z_t(s)=\frac{O_t(x_t(s))-a_t(s)}{b_t(s)}.\] This notation does not select a current, height, or particle-position convention, or assert a common scaling across all initial conditions. A one-point limiting law concerns $Z_t(s)$ for a fixed parameter, whereas a process limit also requires joint distributions and an appropriate mode or topology of convergence. The supplied record does not enumerate the initial-condition classes, fix these normalizations, or specify that topology. The 2026-08-21 triage labels the record partially solved and reports proved Airy-process and Tracy--Widom scaling limits for step and several other initial conditions, while saying a complete long-time classification across initial conditions and models remains open. The displayed source question itself concerns TASEP; the broader mention of models is retained as a literature qualification rather than a new universal statement.
\subsection{Short English statement}
Classify the long-time one-point fluctuation laws and fluctuation processes of TASEP for different initial conditions. The chosen observable, centering, spatial and fluctuation scales, initial-condition classes, and convergence topology must be specified for each regime.
\subsection{Sources}
\begin{itemize}\item Local UnsolvedMath registry, permanent ID 30000839, source problem catalogue entry, \emph{Fluctuation Limits for Exclusion Processes}, proposed year 2007; no individual proposer supplied.\item Supplied citation: \emph{Large Scale Stochastic Dynamics} (2008), Oberwolfach Reports, DOI 10.4171/owr/2007/42, pp.\ 2443--2530: \url{https://doi.org/10.4171/owr/2007/42}. Supplied literature label: original Oberwolfach problem.\item Supplied literature label: TASEP Airy limits, \url{https://link.springer.com/article/10.1007/s00440-021-01074-0}.\item Supplied literature label: Airy-one TASEP limit, \url{https://arxiv.org/abs/math-ph/0611071}.\end{itemize}
\end{document}
